% !TEX root = ../discretemath.tex

\section{Дискретная вероятность и условная вероятность}

\subsection{Конечное вероятностное пространство}

\begin{Definition}{Конечное вероятностное пространство}{}
	Пусть \(U\) — конечное множество элементарных исходов.
	Функция
	\[
	P_r \colon U \to [0,1]
	\]
	называется \textit{вероятностным распределением} на \(U\), если
	\[
	\sum_{\omega \in U} P_r(\omega)=1.
	\]
	Тогда пара \((U, P_r)\) называется \textit{конечным вероятностным пространством}.
\end{Definition}

\begin{Definition}{Событие и его вероятность}{}
	Пусть \((U, P_r)\) — конечное вероятностное пространство.
	Любое подмножество \(A \subseteq U\) называется \textit{событием}.
	Вероятность события \(A\):
	\[
	P_r(A)=\sum_{\omega\in A} P_r(\omega).
	\]
\end{Definition}

\begin{Remark}{}{}
	В случае \emph{равновероятных} исходов ($P_r(\omega)=1/|U|$) вероятность события сводится к классической формуле
	\[
	P_r(A)=\frac{|A|}{|U|}
	\]
	— «число благоприятных исходов, делённое на общее число исходов».
\end{Remark}

\begin{Statement}{Свойства вероятности}{}
	Пусть \(A,B \subseteq U\). Тогда:
	\begin{enumerate}[label=\arabic*.]
		\item \(P_r(\varnothing)=0,\qquad P_r(U)=1.\)
		\item \(P_r(A)+P_r(U\setminus A)=1.\)
		\item \(P_r(A\cup B)=P_r(A)+P_r(B)-P_r(A\cap B).\)
	\end{enumerate}
\end{Statement}

\begin{figure}[H]
	\centering
	\begin{tikzpicture}[font=\small, >=Stealth]
		\begin{scope}
			\clip (0,0) circle (1.1);
			\fill[black!18] (1.1,0) circle (1.1);
		\end{scope}
		\draw[thick] (0,0) circle (1.1);
		\draw[thick] (1.1,0) circle (1.1);
		\node at (-0.5,0) {$A$};
		\node at (1.6,0) {$B$};
		\node at (0.55,0) {\scriptsize$A\cap B$};
		\node[align=center] at (0.55,-1.7)
		{$P_r(A\cup B)=P_r(A)+P_r(B)-P_r(A\cap B)$};
	\end{tikzpicture}
	\caption{Свойство~3: при сложении $P_r(A)+P_r(B)$ пересечение $A\cap B$ учитывается дважды, поэтому его вычитают один раз. Это частный случай формулы включений--исключений для двух событий.}
\end{figure}

\subsection{Примеры}

\begin{Example}{$n$ подбрасываний монеты}{}
	Рассмотрим \(n\) подбрасываний честной монеты. Каждый исход — последовательность из \(n\) символов ($1$ — орёл, $0$ — решка). Найдём вероятность того, что выпало ровно \(k\) орлов.
\end{Example}
Множество исходов:
\[
U=\{0,1\}^n,
\qquad |U|=2^n,
\qquad P_r(\omega)=2^{-n}.
\]
Пусть \(A_k=\{\omega\in U \mid |\omega|=k\}\), где \(|\omega|\) — число орлов. Так как таких последовательностей ровно \(\binom{n}{k}\),
\[
P_r(\text{ровно }k\text{ орлов})
=\binom{n}{k}2^{-n}.
\]

\begin{Example}{7 карт: ровно 5 карт одной масти}{}
	Из стандартной колоды в \(52\) карты случайно выбираются \(7\) карт. Найти вероятность того, что среди них есть \textit{ровно} 5 карт одной масти.
\end{Example}
Выберем масть ($4$ способа). Тогда $5$ карт этой масти выбираются из $13$ — это $\binom{13}{5}$, а оставшиеся $2$ карты — из $39$ карт других мастей — это $\binom{39}{2}$. Общее число наборов $\binom{52}{7}$. Следовательно,
\[
P(\text{ровно 5 карт одной масти})
=\frac{4\binom{13}{5}\binom{39}{2}}{\binom{52}{7}}.
\]

\begin{Example}{Задача о потерянном билете}{}
	В автобус заходит \(N\) человек. У первого пассажира билет потерян.
	\begin{itemize}
		\item Первый пассажир садится на \textit{случайное} место.
		\item Каждый следующий пассажир садится на своё место, если оно свободно; иначе — на случайное из оставшихся свободных мест.
	\end{itemize}
	Найти вероятность того, что \(N\)-й пассажир сядет на своё место.
\end{Example}

\subsubsection*{Решение}

Обозначим через \(a_n\) вероятность того, что \(n\)-й пассажир в итоге сядет на своё место. После первого шага первый пассажир случайно выбирает одно из \(n\) мест. Разберём случаи:
\begin{itemize}
	\item Сел на место \(1\) — все остальные сядут на свои места, вклад \(\tfrac1n\cdot 1.\)
	\item Сел на место \(n\) — \(n\)-й пассажир уже не сможет сесть на своё, вклад \(\tfrac1n\cdot 0.\)
	\item Сел на место \(k\), где \(2\le k\le n-1\) — пассажиры \(2,\dots,k-1\) сядут на свои места, а пассажир \(k\) оказывается в той же задаче для \(n-k+1\) пассажиров, вклад \(\tfrac1n\cdot a_{n-k+1}.\)
\end{itemize}
Следовательно,
\[
a_n=\frac1n\left(1+a_{n-1}+a_{n-2}+\dots+a_2+0\right)
=\frac1n\left(1+\sum_{j=2}^{n-1}a_j\right).
\]

\subsubsection*{База}

При \(n=2\): \(a_2=\tfrac12\) (первый пассажир с равной вероятностью садится на место \(1\) или \(2\)). При \(n=3\) аналогично \(a_3=\tfrac12\).

\begin{Proposition}{}{}
	Для всех \(n\ge 2\) выполнено \(a_n=\dfrac12.\)
\end{Proposition}

\begin{proof}
	Индукция. \textbf{База:} \(a_2=\tfrac12\).

	\textbf{Переход:} пусть \(a_2=\dots=a_{n-1}=\tfrac12\). Тогда
	\[
	a_n
	=\frac1n\left(1+\sum_{j=2}^{n-1}a_j\right)
	=\frac1n\left(1+\frac{n-2}{2}\right)
	=\frac1n\cdot\frac{n}{2}
	=\frac12,
	\]
	поскольку слагаемых в сумме \((n-2)\). Значит \(a_n=\tfrac12\) для всех \(n\ge 2\).
\end{proof}

\begin{Example}{Парадокс дней рождения}{}
	В группе $40$ человек: какова вероятность, что у двух совпадает день рождения?
\end{Example}
Удобнее рассмотреть противоположное событие — \textit{у всех дни рождения различны}. Общее число способов назначить дни рождения: \(365^{40}\). Число благоприятных исходов:
\[
365\cdot 364\cdots 326=\frac{365!}{325!}.
\]
Оценим сверху: среднее первого и последнего множителей равно $\tfrac{365+326}{2}=345.5$, а произведение не превосходит степени среднего арифметического ($ab\le\bigl(\tfrac{a+b}{2}\bigr)^2$), поэтому
\[
P(\text{все различны})
=\frac{365\cdot 364\cdots 326}{365^{40}}
<\left(\frac{345.5}{365}\right)^{40}
=\left(1-\frac{19.5}{365}\right)^{40}.
\]
Следовательно,
\[
P(\text{есть совпадение})
=1-\frac{365\cdot 364\cdots 326}{365^{40}}
\]
оказывается близка к $1$ (уже для $40$ человек — около $0.89$).

\subsection{Формула включений--исключений}

\begin{Theorem}{Формула включений--исключений}{}
	Пусть \(A_1,\dots,A_k\subseteq U\). Тогда
	\[
	P_r\!\left(\bigcup_{i=1}^k A_i\right)
	=
	\sum_{\varnothing\neq I\subseteq \{1,\dots,k\}}
	(-1)^{|I|-1}
	P_r\!\left(\bigcap_{i\in I} A_i\right).
	\]
\end{Theorem}

В развёрнутом виде:
\[
P_r(A_1\cup\cdots\cup A_k)
=\sum_{i=1}^k P_r(A_i)
-\sum_{1\le i<j\le k}P_r(A_i\cap A_j)
+\cdots
+(-1)^{k-1}P_r(A_1\cap\cdots\cap A_k).
\]

\begin{Remark}{}{}
	Свойство~3 из списка свойств вероятности — это случай $k=2$. Важное применение формулы — подсчёт вероятности того, что случайная перестановка не имеет неподвижных точек (см. соответствующий билет о неподвижных точках перестановки).
\end{Remark}

\subsection{Условная вероятность}

\begin{Definition}{Условная вероятность}{}
	Пусть \(A,B\subseteq U\) и \(P_r(B)>0\). \textit{Условной вероятностью события \(A\) при условии \(B\)} называется
	\[
	P_r(A\mid B)=\frac{P_r(A\cap B)}{P_r(B)}.
	\]
\end{Definition}

\begin{Remark}{}{}
	Содержательно: мы «сужаем» вероятностное пространство до $B$ (теперь известно, что $B$ произошло) и нормируем вероятности на $P_r(B)$.
\end{Remark}

\begin{Statement}{}{}
	Пусть \(P_r(C)>0\). Тогда для любых событий \(A,B\)
	\[
	P_r(A\cup B\mid C)
	=
	P_r(A\mid C)+P_r(B\mid C)-P_r(A\cap B\mid C).
	\]
\end{Statement}

\begin{proof}
	По определению $P_r(A\cup B\mid C)=\dfrac{P_r((A\cup B)\cap C)}{P_r(C)}$. Так как
	\[
	(A\cup B)\cap C=(A\cap C)\cup(B\cap C),
	\]
	по свойству~3 числитель равен $P_r(A\cap C)+P_r(B\cap C)-P_r(A\cap B\cap C)$. Деля на $P_r(C)$, получаем требуемое.
\end{proof}

\begin{Statement}{}{}
	Пусть \(P_r(B\cap C)>0\). Тогда
	\[
	P_r(A\mid B\cap C)=\frac{P_r(A\cap B\mid C)}{P_r(B\mid C)}.
	\]
\end{Statement}

\begin{proof}
	По определению $P_r(A\mid B\cap C)=\dfrac{P_r(A\cap B\cap C)}{P_r(B\cap C)}$. С другой стороны,
	\[
	\frac{P_r(A\cap B\mid C)}{P_r(B\mid C)}
	=\frac{P_r(A\cap B\cap C)/P_r(C)}{P_r(B\cap C)/P_r(C)}
	=\frac{P_r(A\cap B\cap C)}{P_r(B\cap C)}
	=P_r(A\mid B\cap C).
	\]
\end{proof}

\begin{Example}{Две монеты}{}
	Два подбрасывания честной монеты. Пусть $A=\{\text{два орла}\}$, $B=\{\text{на первой монете орёл}\}$. Тогда $P_r(A)=\tfrac14$, $P_r(B)=\tfrac12$. Поскольку $A\subseteq B$, имеем $A\cap B=A$, и
	\[
	P_r(A\mid B)=\frac{P_r(A\cap B)}{P_r(B)}=\frac{1/4}{1/2}=\frac12.
	\]
	(Зная, что первая монета — орёл, осталось дождаться орла на второй: вероятность $\tfrac12$.)
\end{Example}

\subsection{Независимость событий}

\begin{Definition}{Независимость двух событий}{}
	События \(A\) и \(B\) называются \textit{независимыми}, если
	\[
	P_r(A\cap B)=P_r(A)\,P_r(B).
	\]
\end{Definition}

\begin{Remark}{}{}
	Если $P_r(B)>0$, это равносильно $P_r(A\mid B)=P_r(A)$: знание $B$ не меняет вероятность $A$. Если $P_r(B)=0$, то $A$ и $B$ независимы автоматически, поскольку $P_r(A\cap B)=0=P_r(A)P_r(B)$.
\end{Remark}

\begin{Example}{Кубик}{}
	Бросается честный кубик, $U=\{1,2,3,4,5,6\}$. Пусть $A=\{\ge 4\}=\{4,5,6\}$, $B=\{\text{делится на }3\}=\{3,6\}$. Тогда $P_r(A)=\tfrac12$, $P_r(B)=\tfrac13$, $A\cap B=\{6\}$, $P_r(A\cap B)=\tfrac16$. Так как $\tfrac12\cdot\tfrac13=\tfrac16$, события \(A\) и \(B\) независимы.
\end{Example}

\subsection{Независимость в совокупности}

\begin{Definition}{Независимость события от семейства событий}{}
	Событие \(A\) \textit{независимо от событий \(B_1,\dots,B_n\)}, если для любого $I\subseteq \{1,\dots,n\}$ события $A$ и $\bigcap_{i\in I} B_i$ независимы.
\end{Definition}

\begin{Definition}{Независимость в совокупности}{}
	События \(A_1,\dots,A_n\) называются \textit{независимыми в совокупности}, если для любого $I\subseteq \{1,\dots,n\}$
	\[
	P_r\!\left(\bigcap_{i\in I} A_i\right)=\prod_{i\in I} P_r(A_i).
	\]
\end{Definition}

\begin{Remark}{}{}
	Попарной независимости \emph{недостаточно} для независимости в совокупности — следующий пример это показывает.
\end{Remark}

\begin{Example}{Попарная независимость без независимости в совокупности}{}
	Рассмотрим \(n\) независимых подбрасываний честной монеты. Пусть $B_i=\{\text{на }i\text{-м подбрасывании орёл}\}$ и $B_{n+1}=\{\text{общее число орлов чётно}\}$. Любые два события из набора $B_1,\dots,B_n,B_{n+1}$ независимы. Однако все вместе они \emph{не} независимы в совокупности: $B_{n+1}$ однозначно определяется по $B_1,\dots,B_n$ (чётность суммы).
\end{Example}

\begin{Statement}{}{}
	Любые \(n\) событий из набора $B_1,\dots,B_n,B_{n+1}$ независимы в совокупности.
\end{Statement}

\begin{proof}
	Два случая.

	\textbf{1) Выбраны только события из \(B_1,\dots,B_n\).} Это независимость результатов отдельных подбрасываний монеты.

	\textbf{2) Среди выбранных есть \(B_{n+1}\), но отсутствует одно из \(B_1,\dots,B_n\).} Без ограничения общности отсутствует $B_n$; рассматриваем $B_1,\dots,B_{n-1},B_{n+1}$. Зафиксируем любые значения первых $n-1$ монет. Тогда исход $n$-й монеты остаётся случайным и равновероятным, и ровно один из двух его исходов делает общее число орлов чётным. Поэтому
	\[
	P_r(B_{n+1}\mid \text{фикс. }B_1,\dots,B_{n-1})=\frac12=P_r(B_{n+1}),
	\]
	то есть $B_{n+1}$ независимо от любой комбинации $B_1,\dots,B_{n-1}$. Значит, эти $n$ событий независимы в совокупности. Любой другой набор из $n$ событий разбирается так же.
\end{proof}

\subsection{Формула Байеса}

\begin{Statement}{Формула Байеса}{}
	Пусть \(P_r(A)>0\) и \(P_r(B)>0\). Тогда
	\[
	P_r(A\mid B)=\frac{P_r(B\mid A)\,P_r(A)}{P_r(B)}.
	\]
\end{Statement}

\begin{proof}
	По определению $P_r(A\mid B)=\dfrac{P_r(A\cap B)}{P_r(B)}$ и $P_r(B\mid A)=\dfrac{P_r(A\cap B)}{P_r(A)}$. Из второго равенства $P_r(A\cap B)=P_r(B\mid A)\,P_r(A)$; подставляя в первое, получаем формулу.
\end{proof}

\begin{figure}[H]
	\centering
	\begin{tikzpicture}[
		>=Stealth, font=\small, thick,
		level 1/.style={sibling distance=52mm, level distance=22mm},
		level 2/.style={sibling distance=24mm, level distance=26mm},
		v/.style={circle, draw, fill=black!80, inner sep=1.6pt},
		el/.style={font=\scriptsize}
		]
		\node[v] {}
		child {node[v] {}
			child {node[v] {} edge from parent node[el,left,pos=0.45]{$P_r(B\mid A)$}}
			child {node[v] {} edge from parent node[el,right,pos=0.45]{$P_r(\bar B\mid A)$}}
			edge from parent node[el,above left]{$P_r(A)$}
		}
		child {node[v] {}
			child {node[v] {} edge from parent node[el,left,pos=0.45]{$P_r(B\mid \bar A)$}}
			child {node[v] {} edge from parent node[el,right,pos=0.45]{$P_r(\bar B\mid \bar A)$}}
			edge from parent node[el,above right]{$P_r(\bar A)$}
		};
	\end{tikzpicture}
	\caption{Дерево вероятностей. Вероятность пути равна произведению вероятностей рёбер: $P_r(A\cap B)=P_r(A)\,P_r(B\mid A)$. Формула Байеса «обращает» условие, выражая $P_r(A\mid B)$ через $P_r(B\mid A)$.}
\end{figure}
